\begin{afterword}
\summaryhead

The post takes up the Mind Projection Fallacy, named by E.~T. Jaynes, who held that probabilities describe partial information in a mind, not properties of things. Its first example is a coin known to be biased in an unknown direction. Before looking at a flip, the frequentist refuses to call heads 0.5, since the coin's long-run frequency is not 0.5. The Bayesian answers that the coin has already landed, that a probability describes what one knows, and that consistent planning requires weighing heads and tails equally. The author sides with the Bayesians: even a fair coin may be predictable by someone who knows enough, as a robot watching an automatic flipper might be, so there is no ``real probability.''

Two puzzles follow. Asking ``Is at least one of your children a boy?'' gives 1/3 for two boys, while asking about the eldest gives 1/2; with four cards, ``at least one ace'' gives 1/5 for a pair of aces, while ``the ace of spades'' gives 1/3. All the numbers are right, the post says, and the paradox disappears once probabilities are seen as states of information: different questions give different evidence.

\responsehead

The mathematics in the post is correct throughout, and the puzzles are stated with the care they need. Enumerating the cases confirms every figure, and the boy-girl puzzle specifies that you ask the question, which is what makes 1/3 the right answer. The hedged claim that a coin's 50 per cent ``may be just plain wrong'' is supported by work published the year before: a coin-tossing machine can be tuned to land heads every time, and natural flips favour the starting side slightly.

The case against the rival views is weaker than the post presents it, though the post warns that it simplifies. Its frequentist merges two positions, frequency and propensity, and says nothing false: this coin's long-run frequency is not 0.5. The Bayesian speaks of how to weigh one flip, a question on which the frequentist offers no rival number, and grants that ``You can call that number whatever you like.'' As the post presents it, the dispute is only about a word. The robot example does raise a real difficulty, the reference class problem, but frequentists and propensity theorists have long treated probability as relative to a class of trials or a set-up, and the post does not give their answer.

The puzzles do not decide between the views. A frequentist counting repeated draws gets 1/5 and 1/3 by the same arithmetic. What the puzzles defeat is the idea that each set of cards carries its own pair-probability, as each apple carries a weight, and frequentists do not hold that. They show that you must condition on exactly what you learned. The post links them to the frequentist through the word ``inherent.''

Two smaller points. The draft-lottery anecdote, offered with ``I believe,'' leaves out that the complaint about mixing was confirmed by the results. And the conclusion is stated for all probabilities, while every example is a classical system whose uncertainty is plainly ignorance. The case that led Popper to propensities, quantum chance, is not discussed.

In short: correct mathematics and well-posed puzzles, together with an argument against frequentism that presents the dispute as one about a word and uses puzzles that frequentists solve the same way.
\end{afterword}
